\documentclass[12pt, a4paper]{article}
\usepackage[utf8]{inputenc}
\usepackage[T1]{fontenc}
\usepackage{lmodern}
\usepackage[margin=1in]{geometry}
\usepackage{setspace}
\usepackage{graphicx}
\usepackage{booktabs}
\usepackage{natbib}
\usepackage{hyperref}
\usepackage{enumitem}

\doublespacing

\title{Gasoline Price Thresholds and Emergency Department Utilization: A Review of the Evidence on Fuel Cost as a Structural Barrier to Acute Care Access}

\author{Ryan Peeler}

\date{}

\begin{document}

\maketitle

\begin{abstract}
Access to emergency medical care in the United States is shaped not only by insurance status and geographic proximity but by the everyday transportation costs patients incur to reach care facilities. This review examines the evidence linking gasoline prices to emergency department (ED) utilization, with particular attention to whether identifiable price thresholds exist beyond which ED visit volumes decline measurably. Drawing on literature from health economics, transportation geography, rural health services research, and behavioral economics, the review synthesizes evidence across three analytic domains: the empirical relationship between fuel costs and ED visit rates, the populations and geographic contexts most sensitive to fuel price fluctuations, and the mechanisms through which gasoline prices influence care-seeking behavior. The evidence indicates that fuel prices exert a measurable effect on ED utilization, but the relationship is non-linear and population-dependent. Rural populations, low-income patients, uninsured individuals, and patients with non-urgent presentations show the greatest sensitivity to fuel cost increases, while patients with high-acuity emergencies show minimal price responsiveness. Several studies identify inflection points in the \$3.50--\$4.50 per gallon range (in 2010--2015 dollars) beyond which ED visits decline for price-sensitive populations, though these thresholds vary by region, population density, and patient demographics. The review argues that gasoline prices function as a regressive structural barrier to emergency care and that health systems and policymakers must account for transportation cost as a social determinant of acute care access.
\end{abstract}

\noindent\textbf{Keywords:} gasoline prices, emergency department utilization, transportation barriers, healthcare access, rural health, health equity, fuel cost threshold, demand elasticity

\section{Introduction}

Emergency departments occupy a unique position in the American healthcare system: they are the only point of care legally required to evaluate and stabilize all patients regardless of ability to pay, under the Emergency Medical Treatment and Labor Act (EMTALA) of 1986. Yet the legal right to emergency care does not guarantee the practical ability to reach it. Among the practical barriers to ED access, transportation cost has received comparatively little scholarly attention.

Gasoline prices in the United States have exhibited substantial volatility, ranging from under \$1.50 per gallon in early 2009 to peaks exceeding \$5.00 per gallon in mid-2022. For patients who depend on personal vehicles to reach emergency care---a category encompassing most of the American population, but disproportionately those in rural and suburban settings without public transit---these fluctuations represent a variable cost of healthcare access that operates independently of the healthcare system itself.

The question of whether gasoline prices affect ED utilization is a question about the price elasticity of demand for emergency care. Standard economic theory predicts that demand for truly emergent care should be highly inelastic with respect to transportation cost. However, a substantial proportion of ED visits involve conditions that patients perceive as less than fully emergent---conditions for which the decision to seek ED care involves weighing severity, urgency, available alternatives, and cost. For these presentations, fuel cost may function as a marginal deterrent.

This review examines the evidence bearing on three questions: Do gasoline prices affect ED utilization rates? Are there identifiable price thresholds beyond which ED visits decline measurably? Which populations and geographic contexts are most sensitive to fuel price effects?

\section{Theoretical Framework: Transportation Cost and Healthcare Demand}

The \citet{grossman1972} model of health capital investment provides a foundational framework: individuals invest in health when the expected return exceeds the full cost of the health-producing input, including transportation. For emergency care, the demand function is complicated by the urgency dimension---transportation cost is most likely to influence demand in the deliberative range of care-seeking decisions.

The distance decay function---the observation that healthcare utilization decreases with increasing distance from the point of care---is one of the most robust findings in health services research \citep{arcury2005}. Fuel cost amplifies this effect by increasing the per-mile cost of travel. Behavioral economics offers additional insight: \citeauthor{thaler1999}'s (\citeyear{thaler1999}) mental accounting framework suggests that fuel cost may be categorized as a transportation expense rather than a healthcare expense, producing threshold effects when prices cross salience boundaries related to reference points and budget constraints \citep{kahneman1979}.

\section{Direct Empirical Evidence}

\citet{wolfe2012} analyzed county-level gasoline prices and ED visit rates, finding a significant negative association: a \$1.00 per gallon increase was associated with approximately a 3--5\% reduction in non-urgent ED visits, with minimal effect on high-acuity visits. \citet{caldwell2013} found that rising fuel costs reduced overall ED volume, driven by non-urgent and semi-urgent triage classifications. \citet{murthy2014} confirmed similar patterns for pediatric ED visits.

\citet{hsia2012} established that transportation distance was a significant predictor of ED utilization and that communities with longer travel distances showed greater sensitivity to fuel cost variation. \citet{skinner2016} found that ED volumes for non-emergent conditions declined by 8--12\% during sustained periods of prices above \$3.50 per gallon in rural Appalachian communities. International evidence from \citet{bergstrand2015} found that high fuel prices were associated with delayed presentation for conditions ultimately requiring emergency care.

\section{Transportation Barriers: The Broader Context}

\citet{syed2013}, in a systematic review of 61 studies, found that transportation barriers were associated with missed appointments, delayed care, and poorer health outcomes, disproportionately affecting low-income, elderly, disabled, and rural populations. \citet{wallace2005} estimated that 3.6 million Americans missed or delayed medical care due to transportation barriers annually. Emergency care is uniquely vulnerable to transportation barriers because visits are unscheduled and often occur during hours when public transportation is unavailable \citep{hsiang2019}.

\citet{chou2019} examined substitution patterns during high gasoline prices, finding that some patients shifted to urgent care or retail clinic settings when available. In communities without these alternatives, there was more evidence of foregone care, indicated by subsequent increases in acute admissions for ambulatory care-sensitive conditions.

\section{Vulnerable Populations}

Rural populations face compounding vulnerabilities: longer driving distances, absence of public transit, and lower household incomes \citep{hsia2012, guttmann2017}. A patient living 25 miles from the nearest ED who drives a vehicle averaging 20 miles per gallon faces a round-trip fuel cost of \$6.25 at \$2.50/gallon but \$12.50 at \$5.00/gallon. Low-income and uninsured populations experience fuel price effects more acutely because transportation cost represents a larger share of household budget \citep{skinner2016}. Elderly populations are partially buffered by higher ambulance utilization rates \citep{cdcnchs2020}.

\section{The Threshold Question}

\citet{wolfe2012} found that the fuel price--ED visit relationship was not well-characterized by a linear model and that the effect intensified above approximately \$3.50 per gallon (2010 dollars). \citet{skinner2016} documented a more pronounced threshold between \$3.50 and \$4.00 per gallon in rural communities.

Threshold mechanisms include mental accounting salience shifts \citep{thaler1999, kahneman1979}, binding budget constraints, and institutional triggers in transportation assistance programs. Despite caveats about temporal and geographic specificity, the convergence around the \$3.50--\$4.50 range suggests a zone within which fuel cost transitions from a non-salient background expense to a material factor in ED care-seeking decisions for price-sensitive populations.

\section{Mechanisms and Alternative Explanations}

The direct fuel cost mechanism---higher prices reducing discretionary income and deterring trips---is supported by evidence that effects are proportional to distance, inversely proportional to income, and concentrated in lower-acuity presentations. Alternative explanations include the general economic activity channel \citep{ruhm2000} and the driving reduction mechanism: \citet{grabowski2004} estimated that a sustained \$1.00 increase in gasoline prices was associated with a 2--3\% reduction in traffic fatalities. However, disaggregation by chief complaint shows that the motor vehicle accident mechanism explains only a fraction of the overall effect.

\citet{caldwell2013} found that fuel price--ED volume associations operated with a 1--2 week lag, consistent with deliberative behavioral change rather than contemporaneous injury reduction.

\section{Implications for Health Systems and Policy}

Fuel price data should be incorporated into ED demand forecasting models. Transportation assistance programs should be extended to unscheduled emergency visits through ride-hailing partnerships and emergency transportation vouchers. Telehealth triage could serve as a partial buffer, though broadband access limitations in affected communities create secondary barriers.

Policy responses could include: inclusion of fuel cost indices in health professional shortage area designation formulas; expansion of community paramedicine programs; investment in rural EMS infrastructure; and exploration of fuel price-indexed emergency transportation subsidies for transportation-vulnerable communities.

\section{Limitations and Future Research}

The empirical literature is relatively small, with most studies from the 2008--2016 period. Threshold estimates are sensitive to specification choices and are not fixed across time and geography. The expansion of urgent care centers, telehealth, and freestanding EDs since 2016 may modify the relationship. Electric vehicle adoption will attenuate the mechanism for some populations, though EV adoption is lowest among those most affected by fuel price barriers.

Future priorities include updated analyses with post-2020 price data, studies of fuel price--telehealth interaction effects, natural experiment designs exploiting fuel tax changes, patient-level surveys on transportation cost in ED decisions, and evaluation of transportation assistance interventions for emergency care access.

\section{Conclusion}

Gasoline prices do affect emergency department utilization, but non-linearly and in a population-specific manner. The relationship is strongest among rural, low-income, and uninsured populations, and concentrated in lower-acuity presentations. High-acuity emergencies show minimal price sensitivity. The threshold evidence converges around the \$3.50--\$4.50 per gallon range as a zone where fuel cost becomes material for price-sensitive populations.

Gasoline prices function as a regressive structural barrier to emergency care---invisible in the design of the healthcare system but consequential in its operation. For a patient who cannot afford the fuel to drive to the ED, the practical effect is the same as lacking insurance: the formal right to care does not translate into actual access. Health systems and policymakers should incorporate transportation cost into their understanding of emergency care access and invest in programs that prevent fuel cost from determining who receives emergency care.

\bibliographystyle{apalike}
\bibliography{references}

\end{document}
